% !TeX program = xelatex
% UTF-8; compile twice with XeLaTeX. No external images or data files needed.
\documentclass[a4paper,landscape,10pt]{article}
\usepackage[margin=14mm,headheight=14pt,headsep=5mm,footskip=7mm]{geometry}
\usepackage{fontspec}
\setmainfont{DejaVu Sans}
\newfontfamily\ipafont{DejaVu Sans}
\usepackage[table]{xcolor}
\usepackage{array,longtable,booktabs,fancyhdr}
\definecolor{ink}{HTML}{18384D}
\definecolor{band}{HTML}{E7EFF4}
\pagestyle{fancy}\fancyhf{}
\fancyhead[L]{\small\color{ink}Parlez Vous QSO?! · Eesti / Estonian}
\fancyhead[R]{\small Englische Phrase · Estnische Übersetzung · IPA}
\fancyfoot[L]{\scriptsize Vorlage: Estonian v 1.0 · Credits: Jüri -- ES5JR}
\fancyfoot[R]{\small\thepage}
\setlength{\parindent}{0pt}
\setlength{\tabcolsep}{3mm}
\renewcommand{\arraystretch}{1.25}
\newcommand{\ipa}[1]{{\ipafont[#1]}}
\begin{document}
{\LARGE\bfseries\color{ink}Parlez Vous QSO?!}\hfill{\large Eesti / Estonian}

\vspace{2mm}
{\small QSO-Sprachführer mit estnischer IPA-Lautschrift. Die Lautschrift bezieht sich auf die estnische Spalte.}

\vspace{2mm}
{\footnotesize \textbf{Leseschlüssel:} \ipa{ˈ} Hauptbetonung, \ipa{ˌ} Nebenbetonung, \ipa{ː} langer Laut;
\ipa{ɑ} = a, \ipa{ɤ} = õ, \ipa{æ} = ä, \ipa{ø} = ö, \ipa{y} = ü, \ipa{ŋ} wie ng; \ipa{ʲ} bezeichnet Palatalisierung.
Grundlage: Standardestnisch. Lang und überlang sind hier zusammengefasst; Schrägstriche trennen Alternativen.}

\vspace{1mm}
\fontsize{9}{11.8}\selectfont
\begin{longtable}{@{}>{\raggedright\arraybackslash}p{13mm}>{\raggedright\arraybackslash}p{73mm}>{\raggedright\arraybackslash}p{75mm}>{\raggedright\arraybackslash}p{90mm}@{}}
\toprule
\textbf{Ref.} & \textbf{English phrase} & \textbf{Eesti} & \textbf{Lautschrift (IPA)}\\
\midrule
\endfirsthead
\toprule
\textbf{Ref.} & \textbf{English phrase} & \textbf{Eesti} & \textbf{Lautschrift (IPA)}\\
\midrule
\endhead
\midrule
\multicolumn{4}{r}{\scriptsize Fortsetzung auf der nächsten Seite}\\
\endfoot
\bottomrule
\endlastfoot
\rowcolor{band} 1.0 & \textbf{Opening Greetings} & \multicolumn{2}{l}{\textbf{Tervitused}} \\*
1.1 & I am very pleased to meet you. & Mul on suur rõõm sind kohata & \ipa{mul on suːr rɤːm sint ˈkohɑtɑ} \\
1.2 & Thank you for your call & Tänan kutsumast & \ipa{ˈtænɑn ˈkutsumɑst} \\
1.3 & Good morning & Tere hommikust & \ipa{ˈtere ˈhomːikust} \\
1.31 & Good day & Tere päevast & \ipa{ˈtere ˈpæevɑst} \\
1.32 & Good afternoon & Tere pärastlõunat & \ipa{ˈtere ˈpærɑstˌlɤunɑt} \\
1.33 & Good evening & Tere õhtust & \ipa{ˈtere ˈɤhtust} \\
1.4 & I know only a little Estonian but I enjoy QSOs in Estonian. & Ma räägin ainult õige pisut eesti keelt, kuid mulle meeldib eesti keeles sidet pidada & \ipa{mɑ ˈræːkin ˈɑinult ˈɤike ˈpisut ˈeːsti keːlt kuit ˈmulːe ˈmeːltip ˈeːsti ˈkeːles ˈsitet ˈpitɑtɑ} \\
\rowcolor{band} 2.0 & \textbf{Operator} & \multicolumn{2}{l}{\textbf{Operaator}} \\*
2.1 & My name is Peter. & Minu nimi on Peter & \ipa{ˈminu ˈnimi on ˈpiːtə} \\
2.2 & I have been licensed since 1960. & Ma olen raadioamatöör olnud alates 1960.-st (tuhande üheksasaja kuuekümnendast) aastast & \ipa{mɑ ˈolen ˈrɑːtioɑmɑˌtøːr ˈolnut ˈɑlɑtes ˈtuhɑnte ˈyheksɑˌsɑjɑ ˈkuːeˌkymːnentɑst ˈɑːstɑst} \\
2.3 & I have been licensed for 43 years. & Ma olen raadioamatöör olnud 43 (nelikümmend kolm) aastat. & \ipa{mɑ ˈolen ˈrɑːtioɑmɑˌtøːr ˈolnut ˈneliˌkymːent kolm ˈɑːstɑt} \\
2.4 & I live in Sittingbourne. & Ma elan Sittingbourne’is. & \ipa{mɑ ˈelɑn ˈsɪtɪŋbɔːnis} \\
2.5 & Sittingbourne is a village. & Sittingbourne on küla. & \ipa{ˈsɪtɪŋbɔːn on ˈkylɑ} \\
2.51 & Sittingbourne is a town. & Sittingbourne on linn. & \ipa{ˈsɪtɪŋbɔːn on linʲː} \\
2.52 & Sittingbourne is a city. & Sittingbourne on (suur)linn. & \ipa{ˈsɪtɪŋbɔːn on linʲː / ˈsuːrˌlinʲː} \\
2.6 & Sittingbourne is between London and Dover. & Sittingbourne asub Londoni ja Doveri vahel. & \ipa{ˈsɪtɪŋbɔːn ˈɑsup ˈlontoni jɑ ˈdəʊvəri ˈvɑhel} \\
2.7 & I live in the South of England. & Ma elan Lõuna-Inglismaal. & \ipa{mɑ ˈelɑn ˈlɤunɑ ˈiŋlisˌmɑːl} \\
2.71 & I live in the North of England. & Ma elan Põhja-Inglismaal. & \ipa{mɑ ˈelɑn ˈpɤhjɑ ˈiŋlisˌmɑːl} \\
2.72 & I live in the East of England. & Ma elan Ida-Inglismaal. & \ipa{mɑ ˈelɑn ˈitɑ ˈiŋlisˌmɑːl} \\
2.73 & I live in West of England. & Ma elan Lääne-Inglismaal. & \ipa{mɑ ˈelɑn ˈlæːne ˈiŋlisˌmɑːl} \\
2.74 & I live in South East England. & Ma elan Kagu-Inglismaal. & \ipa{mɑ ˈelɑn ˈkɑku ˈiŋlisˌmɑːl} \\
2.75 & I live in South West England. & Ma elan Edela-Inglismaal. & \ipa{mɑ ˈelɑn ˈetelɑ ˈiŋlisˌmɑːl} \\
2.76 & I live in North East England. & Ma elan Kirde-Inglismaal. & \ipa{mɑ ˈelɑn ˈkirte ˈiŋlisˌmɑːl} \\
2.77 & I live in North West England. & Ma elan Loode-Inglismaal. & \ipa{mɑ ˈelɑn ˈloːte ˈiŋlisˌmɑːl} \\
2.78 & I live in central England. & Ma elan Kesk-Inglismaal. & \ipa{mɑ ˈelɑn ˈkesk ˈiŋlisˌmɑːl} \\
2.8 & I live in the country. & Ma elan maal. & \ipa{mɑ ˈelɑn mɑːl} \\
2.81 & I live by the sea. & Ma elan mere ääres. & \ipa{mɑ ˈelɑn ˈmere ˈæːres} \\
2.82 & I live near mountains. & Ma elan mägede lähedal. & \ipa{mɑ ˈelɑn ˈmækete ˈlæhetɑl} \\
2.9 & Sittingbourne is 50 km from London. & Sittingbourne on 50 (viiekümne) kilomeetri kaugusel Londonist. & \ipa{ˈsɪtɪŋbɔːn on ˈviːeˌkymːne ˈkiloˌmeːtri ˈkɑukusel ˈlontonist} \\
2.11 & I am retired. & Ma olen pensionil. & \ipa{mɑ ˈolen ˈpensʲionil} \\
2.12 & My home call sign is M0ECS. & Mu kodukutsung on M0ECS. & \ipa{mu ˈkotuˌkutsuŋ on em nulː eː tseː es} \\
\rowcolor{band} 3.0 & \textbf{Station} & \multicolumn{2}{l}{\textbf{Jaam}} \\*
3.1 & My radio is an Icom 706. & Mu (raadio)jaam (transiiver) on Icom 706. & \ipa{mu jɑːm on ˈaɪkɒm ˈseitsesɑtɑ kuːs / mu ˈrɑːtioˌjɑːm on ˈaɪkɒm ˈseitsesɑtɑ kuːs / mu trɑnˈsiːver on ˈaɪkɒm ˈseitsesɑtɑ kuːs} \\
3.2 & My station is all homebrew (homemade). & Ma olen kogu jaama ise ehitanud. & \ipa{mɑ ˈolen ˈkoku ˈjɑːmɑ ˈise ˈehitɑnut} \\
3.3 & I am running 5 Watts. & Jaama võimsus on 5 vatti. & \ipa{ˈjɑːmɑ ˈvɤimsus on viːs ˈvɑtːi} \\
3.4 & My antenna is a long wire. & Mu antenniks on pikk traat. & \ipa{mu ɑnˈtenːiks on pikː trɑːt} \\
3.41 & My antenna is a dipole. & Mul on dipoolantenn. & \ipa{mul on tiˈpoːlˌɑntenː} \\
3.42 & My antenna is a G5RV. & Mu antenniks on G5RV. & \ipa{mu ɑnˈtenːiks on keː viːs er veː} \\
3.43 & My antenna is a doublet. & Mul on dublettantenn. & \ipa{mul on tupˈletːˌɑntenː} \\
3.44 & My antenna is a delta loop. & Mu antenniks on Delta-Loop. & \ipa{mu ɑnˈtenːiks on ˈteltɑ luːp} \\
3.45 & My antenna is a vertical. & Mul on vertikaalantenn. & \ipa{mul on vertiˈkɑːlˌɑntenː} \\
3.46 & My antenna is a log periodic. & Mul on logoperioodiline antenn. & \ipa{mul on ˈlokoˌperioːtiline ɑnˈtenː} \\
3.47 & My antenna is a Yagi beam. & Mul on yagiantenn. & \ipa{mul on ˈjɑkiˌɑntenː} \\
3.48 & My antenna is a 5 element beam. & Mul on viie-elemendiline yagi. & \ipa{mul on ˈviːe ˌelemenˈtiline ˈjɑki} \\
\rowcolor{band} 4.0 & \textbf{Conditions} & \multicolumn{2}{l}{\textbf{Olud}} \\*
4.1 & Thank you for my 599 report. & Tänan 599 raporti eest. & \ipa{ˈtænɑn viːs ˈyheksɑ ˈyheksɑ ˈrɑporti eːst} \\
4.2 & The band conditions are good. & Levitingimused on head. & \ipa{ˈleviˌtiŋkimuset on heɑt} \\
4.3 & The band conditions are bad. & Levitingimused on halvad. & \ipa{ˈleviˌtiŋkimuset on ˈhɑlvɑt} \\
4.4 & There is strong interference here. & Siin on tugevad häired. & \ipa{siːn on ˈtukevɑt ˈhæiret} \\
4.5 & Can you change to the 80m band? & Kas sa saaksid minna 80 (kaheksakümnele) meetrile? & \ipa{kɑs sɑ ˈsɑːksit ˈminːɑ ˈkɑheksɑˌkymːnele ˈmeːtrile} \\
4.6 & Please change frequency to 3.558 MHz. & Palun mine sagedusele 3,558 MHz. & \ipa{ˈpɑlun ˈmine ˈsɑketusele kolm ˈkomɑ viːs viːs ˈkɑheksɑ ˈmekɑˌhertsi} \\
\rowcolor{band} 5.0 & \textbf{General Queries and Comments} & \multicolumn{2}{l}{\textbf{Üldised küsimused ja kommentaarid}} \\*
5.1 & What does <\#\#\#\#\#> mean please? & Palun ütle, mida tähendab <\#\#\#\#\#>? & \ipa{ˈpɑlun ˈytle ˈmitɑ ˈtæhentɑp …} \\
5.2 & Please wait a moment, I need to check my Estonian dictionary. & Oota üks hetk palun, ma vaatan eesti keele sõnastikust järele. & \ipa{ˈoːtɑ yks hetk ˈpɑlun mɑ ˈvɑːtɑn ˈeːsti ˈkeːle ˈsɤnɑstikust ˈjærele} \\
5.3 & How do you say, <\#\#\#\#\#> in Estonian ? & Kuidas on <\#\#\#\#\#> eesti keeles? & \ipa{ˈkuitɑs on … ˈeːsti ˈkeːles} \\
5.5 & My favourite modes of operation are CW / RTTY / PSK / MT63 / Olivia / Hellschreiber / Digital Modes / and SSB & Minu lemmiktööliigid on telegraaf / raadioteletaip / PSK / MT63 / Olivia / Hell / digitaalsed tööliigid / ja SSB. & \ipa{ˈminu ˈlemːikˌtøːˌliːkit on teleˈkrɑːf / ˈrɑːtioˌteleˌtɑip / peː es kɑː / em teː ˈkuːsˌkymːent kolm / oˈliviɑ / hɛl / tikiˈtɑːlset ˈtøːˌliːkit / jɑ es es peː} \\
5.6 & My best DX is Australia. & Mu kõige kaugem side on olnud Austraaliaga. & \ipa{mu ˈkɤike ˈkɑukem ˈsite on ˈolnut ɑusˈtrɑːliɑkɑ} \\
5.7 & I enjoy working contests. & Mulle meeldib töötada võistlustel. & \ipa{ˈmulːe ˈmeːltip ˈtøːtɑtɑ ˈvɤistlustel} \\
5.8 & I like working QRP. & Mulle meeldib väikese võimsusega töötamine. & \ipa{ˈmulːe ˈmeːltip ˈvæikese ˈvɤimsusekɑ ˈtøːtɑmine} \\
\rowcolor{band} 6.0 & \textbf{The Weather} & \multicolumn{2}{l}{\textbf{Ilm}} \\*
6.1 & It is sunny. & Ilm on päikeseline. & \ipa{ilm on ˈpæikeseˌline} \\
6.2 & It is windy. & Ilm on tuuline. & \ipa{ilm on ˈtuːline} \\
6.3 & It is raining. & Sajab. & \ipa{ˈsɑjɑp} \\
6.4 & It is snowing. & Sajab lund. & \ipa{ˈsɑjɑp lunt} \\
6.5 & It is very hot. & Ilm on väga kuum. & \ipa{ilm on ˈvækɑ kuːm} \\
6.6 & It is very cold. & Ilm on väga külm. & \ipa{ilm on ˈvækɑ kylm} \\
6.7 & The temperature is 8C. & Temperatuur on 8 (kaheksa) kraadi. & \ipa{ˌtemperɑˈtuːr on ˈkɑheksɑ ˈkrɑːti} \\
6.8 & It is overcast. & Ilm on pilvine. & \ipa{ilm on ˈpilvine} \\
6.9 & It is foggy. & Ilm on udune. & \ipa{ilm on ˈutune} \\
6.11 & It has been raining all day. & Kogu päeva on sadanud. & \ipa{ˈkoku ˈpæevɑ on ˈsɑtɑnut} \\
6.12 & There is 20cm of snow here. & Lume paksus on 20 (kakskümmend) sentimeetrit. & \ipa{ˈlume ˈpɑksus on ˈkɑksˌkymːent ˌsentiˈmeːtrit} \\
\rowcolor{band} 8.0 & \textbf{QSL} & \multicolumn{2}{l}{\textbf{QSL-kaardid}} \\*
8.1 & My QSL card is sure via the bureau & Saadan kaardi büroo kaudu. & \ipa{ˈsɑːtɑn ˈkɑːrti pyˈroː ˈkɑutu} \\
8.2 & QSL via eQSL & Saadan eQSL-i. & \ipa{ˈsɑːtɑn iː kuː es eli} \\
8.3 & No eQSL please & Palun ära saada eQSL-i. & \ipa{ˈpɑlun ˈærɑ ˈsɑːtɑ iː kuː es eli} \\
8.4 & I will send you my card direct. & Saadan sulle oma kaardi otse. & \ipa{ˈsɑːtɑn ˈsulːe ˈomɑ ˈkɑːrti ˈotse} \\
8.5 & What is your address? & Milline on sinu aadress? & \ipa{ˈmilːine on ˈsinu ˈɑːtresː} \\
8.6 & Please send your QSL card direct. & Palun saada oma QSL-kaart otse. & \ipa{ˈpɑlun ˈsɑːtɑ ˈomɑ kuː es el kɑːrt ˈotse} \\
8.7 & My address is... & Minu aadress on… & \ipa{ˈminu ˈɑːtresː on …} \\
8.8 & My address is shown on QSL.com / EHam.net & Minu aadressi saab QSL.com-ist / eHam.net-ist. & \ipa{ˈminu ˈɑːtresːi sɑːp kuː es el puŋkt komist / iː hæm puŋkt netist} \\
8.9 & Sorry I do not send QSL cards. & Palun vabandust, ma ei saada QSL-kaarte. & \ipa{ˈpɑlun ˈvɑpɑntust mɑ ei ˈsɑːtɑ kuː es el ˈkɑːrte} \\
8.11 & Please send QSL card via G1JQH. & Palun saada QSL-kaart G1JQH kaudu. & \ipa{ˈpɑlun ˈsɑːtɑ kuː es el kɑːrt keː yks jotː kuː hɑː ˈkɑutu} \\
8.12 & My WACRAL / GQRP / DXCC / FISTS number is 1234. & Minu WACRAL / GQRP / DXCC / FISTS number on 1234 (tuhat kakssada kolmkümmend neli). & \ipa{ˈminu ˈwakral / keː kuː er peː / teː iks tseː tseː / fɪsts ˈnumper on ˈtuhɑt ˈkɑksˌsɑtɑ ˈkolmˌkymːent ˈneli} \\
\rowcolor{band} 9.0 & \textbf{Final Greetings} & \multicolumn{2}{l}{\textbf{Lõputervitused}} \\*
9.1 & Thank you. & Tänan sind. & \ipa{ˈtænɑn sint} \\
9.2 & Thank you for a very enjoyable contact. & Tänan sind meeldiva side eest. & \ipa{ˈtænɑn sint ˈmeːltivɑ ˈsite eːst} \\
9.3 & I enjoyed our chat. & Oli tore vestlus (jutuajamine). & \ipa{ˈoli ˈtore ˈvestlus / ˈoli ˈtore ˈjutuˌɑjɑmine} \\
9.4 & I enjoyed our QSO in Estonian. & Mulle meeldis väga meie eestikeelne side. & \ipa{ˈmulːe ˈmeːltis ˈvækɑ ˈmeie ˈeːstiˌkeːlne ˈsite} \\
9.5 & It has been great fun. & Oli väga lõbus (tore). & \ipa{ˈoli ˈvækɑ ˈlɤpus / ˈoli ˈvækɑ ˈtore} \\
9.6 & If you would like a regular contact with me, you can contact me via email. & Kui tahad minuga pidevalt suhelda, siis saad saata e-kirja. & \ipa{kui ˈtɑhɑt ˈminukɑ ˈpitevɑlt suˈheltɑ siːs sɑːt ˈsɑːtɑ ˈeːˌkirjɑ} \\
9.7 & My email address is... & Mu e-posti aadress on… & \ipa{mu ˈeːˌposti ˈɑːtresː on …} \\
9.8 & My email address is shown on QSL.com / EHam.net & Mu e-posti aadress on kirjas QSL.com-is / EHam.net-is & \ipa{mu ˈeːˌposti ˈɑːtresː on ˈkirjɑs kuː es el puŋkt komis / iː hæm puŋkt netis} \\
9.9 & I look forward to meeting you again. & Loodan sind jälle kohata. & \ipa{ˈloːtɑn sint ˈjælːe ˈkohɑtɑ} \\
9.11 & My best wishes to you and your family. & Minu parimad soovid sulle ja perele. & \ipa{ˈminu ˈpɑrimɑt ˈsoːvit ˈsulːe jɑ ˈperele} \\
9.12 & Best wishes to your wife. & Minu parimad soovid sinu abikaasale (naisele). & \ipa{ˈminu ˈpɑrimɑt ˈsoːvit ˈsinu ˈɑpiˌkɑːsɑle / ˈminu ˈpɑrimɑt ˈsoːvit ˈsinu ˈnɑisele} \\
9.121 & Best wishes to your husband. & Minu parimad soovid sinu abikaasale (mehele) & \ipa{ˈminu ˈpɑrimɑt ˈsoːvit ˈsinu ˈɑpiˌkɑːsɑle / ˈminu ˈpɑrimɑt ˈsoːvit ˈsinu ˈmehele} \\
9.13 & Good luck. & Õnne edasiseks. & \ipa{ˈɤnːe ˈetɑsiseks} \\
9.14 & Good health. & Head tervist. & \ipa{heɑt ˈtervist} \\
9.15 & Goodbye & Head aega. & \ipa{heɑt ˈɑekɑ} \\
9.16 & I wish you an enjoyable afternoon. & Soovin sulle kena pärastlõunat. & \ipa{ˈsoːvin ˈsulːe ˈkenɑ ˈpærɑstˌlɤunɑt} \\
9.161 & I wish you an enjoyable evening. & Soovin sulle kena õhtut. & \ipa{ˈsoːvin ˈsulːe ˈkenɑ ˈɤhtut} \\
9.162 & I wish you an enjoyable weekend. & Soovin sulle kena nädalalõppu. & \ipa{ˈsoːvin ˈsulːe ˈkenɑ ˈnætɑlɑˌlɤpːu} \\
9.17 & Enjoy your holiday. & Ilusat puhkust. & \ipa{ˈilusɑt ˈpuhkust} \\
9.18 & Happy New Year. & Head vana aasta lõppu! (zwischen Weihnachten und 31. Dezember)\newline Head uut aastat! (31. Dezember bis 6. Januar) & \ipa{heɑt ˈvɑnɑ ˈɑːstɑ ˈlɤpːu / heɑt uːt ˈɑːstɑt} \\
\rowcolor{band} 10.0 & \textbf{Country} & \multicolumn{2}{l}{\textbf{Riik}} \\*
10.1 & Albania & Albaania & \ipa{ɑlˈpɑːniɑ} \\
10.2 & Austria & Austria & \ipa{ˈɑustriɑ} \\
10.3 & Belgium & Belgia & \ipa{ˈpelkiɑ} \\
10.4 & Belorussia & Valgevene & \ipa{ˈvɑlkeˌvene} \\
10.5 & Bosnia \& Herzegovina & Bosnia-Hertsegoviina & \ipa{ˈposniɑ hertsekoˈviːnɑ} \\
10.6 & Bulgaria & Bulgaaria & \ipa{pulˈkɑːriɑ} \\
10.7 & Croatia & Horvaatia & \ipa{horˈvɑːtiɑ} \\
10.8 & Czech Republic & Tšehhi (Vabariik) & \ipa{ˈtʃehːi / ˈtʃehːi ˈvɑpɑˌriːk} \\
10.9 & Denmark & Taani & \ipa{ˈtɑːni} \\
10.11 & Eire & Iirimaa & \ipa{ˈiːriˌmɑː} \\
10.12 & England & Inglismaa & \ipa{ˈiŋlisˌmɑː} \\
10.13 & Estonia & Eesti & \ipa{ˈeːsti} \\
10.14 & Finland & Soome & \ipa{ˈsoːme} \\
10.15 & France & Prantsusmaa & \ipa{ˈprɑntsusˌmɑː} \\
10.16 & Germany & Saksamaa & \ipa{ˈsɑksɑˌmɑː} \\
10.17 & Gibraltar & Gibraltar & \ipa{tʃipˈrɑltɑr} \\
10.18 & Greece & Kreeka & \ipa{ˈkreːkɑ} \\
10.19 & Holland & Holland & \ipa{ˈholːɑnt} \\
10.21 & Hungary & Ungari & \ipa{ˈuŋkɑri} \\
10.22 & Iceland & Island & \ipa{ˈislɑnt} \\
10.23 & Northern Ireland & Põhja-Iirimaa & \ipa{ˈpɤhjɑ ˈiːriˌmɑː} \\
10.24 & Isle of Man & Mani saar & \ipa{ˈmɑni sɑːr} \\
10.25 & Italy & Itaalia & \ipa{iˈtɑːliɑ} \\
10.26 & Latvia & Läti & \ipa{ˈlæti} \\
10.27 & Liechtenstein & Liechtenstein & \ipa{ˈlɪçtn̩ʃtaɪn} \\
10.28 & Lithuania & Leedu & \ipa{ˈleːtu} \\
10.29 & Luxembourg & Luksemburg & \ipa{ˈluksemˌpurk} \\
10.31 & Macedonia & Makedoonia & \ipa{mɑkeˈtoːniɑ} \\
10.32 & Moldavia & Moldaavia & \ipa{molˈtɑːviɑ} \\
10.33 & Netherlands & Holland (=Madalmaad) & \ipa{ˈholːɑnt / ˈmɑtɑlˌmɑːt} \\
10.34 & Norway & Norra & \ipa{ˈnorːɑ} \\
10.35 & Poland & Poola & \ipa{ˈpoːlɑ} \\
10.36 & Portugal & Portugal & \ipa{ˈportukɑl} \\
10.37 & Romania & Rumeenia & \ipa{ruˈmeːniɑ} \\
10.38 & Russian Federation & Venemaa (Vene Föderatsioon) & \ipa{ˈveneˌmɑː / ˈvene ˌfɤterɑˈtsioːn} \\
10.39 & Scotland & Šotimaa & \ipa{ˈʃotiˌmɑː} \\
10.41 & Slovakia & Slovakkia & \ipa{sloˈvɑkːiɑ} \\
10.42 & Slovenia & Sloveenia & \ipa{sloˈveːniɑ} \\
10.43 & Spain & Hispaania & \ipa{hisˈpɑːniɑ} \\
10.44 & Svalbard & Svalbard & \ipa{ˈsvɑːlbɑr} \\
10.45 & Sweden & Rootsi & \ipa{ˈroːtsi} \\
10.46 & Switzerland & Šveits & \ipa{ʃveits} \\
10.47 & Turkey & Türgi & \ipa{ˈtyrki} \\
10.48 & Ukraine & Ukraina & \ipa{ukˈrɑinɑ} \\
10.49 & Wales & Wales & \ipa{weɪlz} \\
\end{longtable}
\normalsize
\clearpage
\section*{Hinweise zur Übertragung und Aussprache}
\fontsize{8.5}{10}\selectfont\raggedright
\textbf{Vorlage und Umfang.} Grundlage ist „Estonian v 1.0.doc“, Credit: Jüri, ES5JR.
Alle 142 Phrasen und Ländernamen sind enthalten. Die ursprünglichen Nummernsprünge bleiben erhalten,
einschließlich des fehlenden Abschnitts 7. Referenz „132“ wurde zu „1.32“ berichtigt.

\smallskip
\textbf{Korrekturen.} G5GV wurde gemäß der englischen Spalte zu G5RV (3.42).
Die beiden stehen gebliebenen Verweise auf German in der englischen Spalte wurden zu Estonian (5.3/9.4).
Doppelte Punkte und Leerzeichen vor Satzzeichen wurden bereinigt; offene Angaben stehen als Auslassung.
Die englischen Datumserläuterungen bei den Neujahrsgrüßen in 9.18 wurden ins Deutsche übertragen;
die estnischen Wünsche bleiben unverändert.

\smallskip
\textbf{Quellenbesonderheiten.} Die Vorlage erklärt mehrere Zahlen bereits in Klammern;
in der IPA-Spalte wird die Zahl jeweils nur einmal gesprochen, in der zum Satz passenden Form.
„(suur)linn“ in 2.52 erscheint als linn beziehungsweise suurlinn.
Die Varianten raadiojaam, jaam und transiiver in 3.1 sind als vollständige Alternativen wiedergegeben.
Die Beispiele Peter, Sittingbourne, 1960 und 43 Jahre sowie historische Ländernamen und Webadressen bleiben erhalten.

\smallskip
\textbf{Bedeutungsunterschiede.} 2.2/2.3 beschreiben die Tätigkeit als Funkamateur statt ausdrücklich eine Lizenz.
5.5 verwendet telegraaf, raadioteletaip und Hell statt wörtlich CW, RTTY und Hellschreiber;
5.8 beschreibt das Arbeiten mit geringer Leistung ohne das Kürzel QRP.
6.7 nennt acht Grad ohne Celsius. 8.1 sagt, dass eine Karte über das Büro geschickt wird,
ohne das englische „sure“ ausdrücklich wiederzugeben.
Die IPA-Spalte folgt jeweils der estnischen Formulierung.

\smallskip
\textbf{Aussprachebasis.} Die manuelle IPA-Spalte ist eine breite Lesehilfe für Standardestnisch.
\ipa{ɑ} steht für a, \ipa{ɤ} für das ungerundete õ, \ipa{æ} für ä, \ipa{ø} für ö und \ipa{y} für ü.
Diphthonge stehen hier als Vokalfolgen; j und gelegentliche Übergangslaute werden vereinfacht dargestellt.
\ipa{ˈ} bezeichnet die Hauptbetonung, \ipa{ˌ} ausgewählte Nebenbetonungen, besonders in Zusammensetzungen.
Viele einheimische Wörter tragen die Betonung am Anfang; Fremdwörter können davon abweichen.

\smallskip
\textbf{Längen und Quantität.} Estnisch unterscheidet kurze, lange und überlange Quantität,
die sich aus Lautdauer, Silbenverhältnissen und Tonverlauf ergibt.
Hier werden lange und überlange Laute gemeinsam mit \ipa{ː} dargestellt.
Die zweite und dritte Quantitätsstufe sind somit nicht vollständig unterschieden;
halblange unbetonte Vokale und genaue Tonverläufe sind ebenfalls nicht kodiert.
Das ist eine bewusste Vereinfachung dieser Aussprachehilfe, keine Behauptung, dass die beiden Stufen gleich wären.

\smallskip
\textbf{Konsonanten.} Geschriebenes b, d und g werden hier überwiegend als stimmlose \ipa{p, t, k}
wiedergegeben; im flüssigen Gespräch können sie teilweise stimmhaft werden.
Lange Konsonanten stehen etwa als \ipa{nː, lː, tː}; \ipa{ŋ} ist die Nasalvariante vor k/g.
\ipa{ʲ} markiert ausgewählte Palatalisierung, besonders in linn und pensionil.
Nicht jede feine Palatalisierung oder Assimilation ist eingetragen.
H wird für deutliches Vorlesen dargestellt; im Gespräch kann es am Wortanfang entfallen.
R steht vereinfacht als \ipa{r}.

\smallskip
\textbf{Namen, Kürzel und Adressen.} Peter, Sittingbourne, Dover und Icom verwenden englische Lesarten,
mit estnischen Endungen, wo die Vorlage sie verlangt.
Delta-Loop wird mit estnischem delta und englischem Loop gelesen; yagi steht als jagi.
Hell und Liechtenstein sind deutsch, Svalbard norwegisch, Wales englisch gelesen.
Andere Ländernamen folgen der estnischen Form. Rufzeichen und Kürzel verwenden estnische Buchstabennamen:
Q „kuu“, S „ess“, L „ell“, G „gee“, R „err“, C „tsee“, J „jott“, H „haa“.
WACRAL und FISTS werden als Wörter gesprochen; das e in eQSL/eHam als englisches \ipa{iː}.
Webpunkte sind „punkt“; die angehängten estnischen Endungen -ist/-is werden mitgesprochen.

\smallskip
\textbf{Zahlen und Klammern.} 1960.-st wird als „tuhande üheksasaja kuuekümnendast“ gelesen,
43 als „nelikümmend kolm“, 50 als „viiekümne“, 706 als „seitsesada kuus“;
599 ziffernweise „viis üheksa üheksa“, MT63 mit „kuuskümmend kolm“.
3,558 MHz = „kolm koma viis viis kaheksa megahertsi“;
1234 = „tuhat kakssada kolmkümmend neli“. Einheiten werden ausgesprochen.
Klammeralternativen in 9.3, 9.5, 9.12, 9.121 und bei Ländern sind als Alternativen enthalten.
Die Datumshinweise in 9.18 werden nicht mitgesprochen. \ipa{…} kennzeichnet offene Angaben und Platzhalter.

\smallskip
\textbf{Fachhinweise.} Eva Liina Asu und Pire Teras: „Estonian“, Journal of the International
Phonetic Association 39(3), 2009, S. 367--372, DOI 10.1017/S002510030999017X.
Eesti Keele Instituut: „Eesti keel IPA transkriptsioonis“, teatmik.eki.ee, 23. April 2023.
Die Konventionen wurden für diese breite, vereinfachte Lesehilfe angepasst.
\end{document}
